% Options for packages loaded elsewhere
\PassOptionsToPackage{unicode}{hyperref}
\PassOptionsToPackage{hyphens}{url}
%
\documentclass[
]{article}
\usepackage{amsmath,amssymb}
\usepackage{iftex}
\ifPDFTeX
  \usepackage[T1]{fontenc}
  \usepackage[utf8]{inputenc}
  \usepackage{textcomp} % provide euro and other symbols
\else % if luatex or xetex
  \usepackage{unicode-math} % this also loads fontspec
  \defaultfontfeatures{Scale=MatchLowercase}
  \defaultfontfeatures[\rmfamily]{Ligatures=TeX,Scale=1}
\fi
\usepackage{lmodern}
\ifPDFTeX\else
  % xetex/luatex font selection
\fi
% Use upquote if available, for straight quotes in verbatim environments
\IfFileExists{upquote.sty}{\usepackage{upquote}}{}
\IfFileExists{microtype.sty}{% use microtype if available
  \usepackage[]{microtype}
  \UseMicrotypeSet[protrusion]{basicmath} % disable protrusion for tt fonts
}{}
\makeatletter
\@ifundefined{KOMAClassName}{% if non-KOMA class
  \IfFileExists{parskip.sty}{%
    \usepackage{parskip}
  }{% else
    \setlength{\parindent}{0pt}
    \setlength{\parskip}{6pt plus 2pt minus 1pt}}
}{% if KOMA class
  \KOMAoptions{parskip=half}}
\makeatother
\usepackage{xcolor}
\usepackage[margin=1in]{geometry}
\usepackage{longtable,booktabs,array}
\usepackage{calc} % for calculating minipage widths
% Correct order of tables after \paragraph or \subparagraph
\usepackage{etoolbox}
\makeatletter
\patchcmd\longtable{\par}{\if@noskipsec\mbox{}\fi\par}{}{}
\makeatother
% Allow footnotes in longtable head/foot
\IfFileExists{footnotehyper.sty}{\usepackage{footnotehyper}}{\usepackage{footnote}}
\makesavenoteenv{longtable}
\usepackage{graphicx}
\makeatletter
\def\maxwidth{\ifdim\Gin@nat@width>\linewidth\linewidth\else\Gin@nat@width\fi}
\def\maxheight{\ifdim\Gin@nat@height>\textheight\textheight\else\Gin@nat@height\fi}
\makeatother
% Scale images if necessary, so that they will not overflow the page
% margins by default, and it is still possible to overwrite the defaults
% using explicit options in \includegraphics[width, height, ...]{}
\setkeys{Gin}{width=\maxwidth,height=\maxheight,keepaspectratio}
% Set default figure placement to htbp
\makeatletter
\def\fps@figure{htbp}
\makeatother
\setlength{\emergencystretch}{3em} % prevent overfull lines
\providecommand{\tightlist}{%
  \setlength{\itemsep}{0pt}\setlength{\parskip}{0pt}}
\setcounter{secnumdepth}{-\maxdimen} % remove section numbering
\ifLuaTeX
  \usepackage{selnolig}  % disable illegal ligatures
\fi
\IfFileExists{bookmark.sty}{\usepackage{bookmark}}{\usepackage{hyperref}}
\IfFileExists{xurl.sty}{\usepackage{xurl}}{} % add URL line breaks if available
\urlstyle{same}
\hypersetup{
  pdftitle={MATH5090M Coursework},
  pdfauthor={Athang Ghag (201712928)},
  hidelinks,
  pdfcreator={LaTeX via pandoc}}

\title{MATH5090M Coursework}
\author{Athang Ghag (201712928)}
\date{30/04/2023}

\begin{document}
\maketitle

\hypertarget{introduction}{%
\subsubsection{Introduction}\label{introduction}}

The information and analysis in this report is based on a dataset of
cancer patients. In this report, we analyze the age at which cancer
patients die and the various factors affecting it. We explore the
relationship between age at which the cancer was diagnosed and the age
at death of patients. We also design a new public health intervention,
which aims to improve care for cancer patients.

\hypertarget{methodology-and-analysis}{%
\subsubsection{Methodology and
Analysis}\label{methodology-and-analysis}}

The given dataset of cancer patients has 210 rows and 18 columns. There
are 210 patients and 18 features with information about each patient.
There are no missing values in the dataset. The features contain
information such as age at which the patient was diagnosed
(\texttt{AgeDiagnosis}) and age at which the patient died
(\texttt{AgeDeath}) which are given in years in numeric format. There is
an unique identity number associated with each patient in the
\texttt{IDnumber} column and similarly a site identity number
(\texttt{SiteID}) for hospitals where the patient is taking treatment.

The data has categorical binary variables such as Sex, Race, Employment
Status, Social Grade Status to give information about patient's
demographics. Moreover, the binary variable \texttt{HighQualityOfLife}
shows whether the patient is leading a high quality life or not.

The categorical variable \texttt{Type} shows the type of cancer the
patients is diagnosed with such as Kidney cancer, Thyroid cancer,
Leukemia, Bowel, Brain, and more. The binary variable
\texttt{Metastases} indicates if the cancer is metastasized or not,
\texttt{Stage3or4} indicates if the cancer has progressed to the third
stage, \texttt{MultiSite} variable indicates if the cancer has spread to
multiple sites, and similarly \texttt{Remission} variable shows if the
cancer is in remission or not. The categorical variable
\texttt{Treatment} shows the type of treatment the patients is
undergoing which include Chemotherapy, Radiotherapy, Immunotherapy, and
No Treatment.

The patients are taking treatment at different hospitals. We group by
the \texttt{SiteID} to calculate number of patients at each hospital.

\begin{longtable}[]{@{}cc@{}}
\caption{Table 1- Number of Patients in each site}\tabularnewline
\toprule\noalign{}
SiteID & Number of Patients \\
\midrule\noalign{}
\endfirsthead
\toprule\noalign{}
SiteID & Number of Patients \\
\midrule\noalign{}
\endhead
\bottomrule\noalign{}
\endlastfoot
1 & 8 \\
5 & 12 \\
6 & 14 \\
11 & 25 \\
13 & 14 \\
16 & 19 \\
17 & 33 \\
19 & 6 \\
20 & 37 \\
21 & 15 \\
22 & 13 \\
25 & 14 \\
\end{longtable}

There are 12 sites in total. Site 20 has the highest number of patients
which is 37 and Site 19 has the lowest number of patients.

We now analyze our target variable \texttt{AgeDeath} and covariate
\texttt{AgeDiagnosis}. We plot the histograms of these two variables to
explore their behaviour.

\begin{figure}

{\centering \includegraphics{assignment1_files/figure-latex/yhist-1} 

}

\caption{Fig 1- Histograms of Patient's Age at Diagnosis and Patients Age at Death.}\label{fig:yhist}
\end{figure}

We observe that, the histogram of patient's age at diagnosis is slightly
right-skewed. The most frequent diagnosis age falls between the ages of
35 to 50 years with the mean age being around 45 years. The frequency of
age at diagnosis above 60 years is low indicating that unfortunately
many people die before they are diagnosed. The red vertical lines in the
graphs show the mean of diagnosis and death ages.

The histogram of patient's age at diagnosis is also slightly
right-skewed but less prominent than the previous histogram. The
majority of patients die between the ages of 45 to 60 years with the
average age being around 54 years. In our data, patients with cancer
don't survive above the age of 80 years. The average age at death is
higher than age at diagnosis as we expect. It is hard to tell with
certainity with histograms if these 2 variable follow normal
distribution so we make use of QQ Plots to assess the normality of our
data.

\begin{figure}

{\centering \includegraphics{assignment1_files/figure-latex/yqq-1} 

}

\caption{Fig 2- QQ Plots of Patient's Age at Diagnosis and Patients Age at Death.}\label{fig:yqq}
\end{figure}

We observe that, for dianosis age plot the majority of points fall on
the normal line but they just curve away at upper and lower ends.
Similarly for death age plot, some points curve away from the normal
line at upper and lower ends. However, these qq plots are close enough
to a normal distribution qq plot. So, for practical purposes we can
assume that both the variables follow normal distribution and perform
statistical tests.

After assuming a normal distribution for our target and dependent
variable we implement a single-level regression model on them. We fit a
single-level regression model using \texttt{lm()} function in R. The
estimated parameters of the regression are as follows:

\begin{verbatim}
Residuals:
    Min      1Q  Median      3Q     Max 
-9.3275 -3.2145 -0.1056  2.4383 13.3094 

Coefficients:
             Estimate Std. Error t value Pr(>|t|)    
(Intercept)  13.51652    2.06384   6.549 4.45e-10 ***
AgeDiagnosis  0.90176    0.04539  19.866  < 2e-16 ***
---
Signif. codes:  0 ‘***’ 0.001 ‘**’ 0.01 ‘*’ 0.05 ‘.’ 0.1 ‘ ’ 1

Residual standard error: 4.632 on 208 degrees of freedom
Multiple R-squared:  0.6549,    Adjusted R-squared:  0.6532 
F-statistic: 394.7 on 1 and 208 DF,  p-value: < 2.2e-16
\end{verbatim}

The p-value of our model is very close to zero so we fail to accept the
null hypothesis which indicates that there is a significant relationship
between age at death and age of diagnosis. The estimated intercept is
13.51 which indicates that on average there is a difference of about 13
years between age at death and age at diagnosis. The estimated slope is
0.9, a positive value but less than 1, which suggests that a diagnosis
of cancer reduces the patient's expected lifespan by 10\%.

\begin{figure}

{\centering \includegraphics{assignment1_files/figure-latex/reg-1} 

}

\caption{Fig 3- A single-level regression model for Age at Death conditional on Age at Diagnosis.}\label{fig:reg}
\end{figure}

The linear regression plot shows a positive relationship between age at
diagnosis and age at death which is also indicated by the estimated
R-squared value of 0.655. The majority of observations are between ages
40 to 55 years. The regression line seems to be a good fit for data
considering majority of the data points are clustered near the line with
the estimated residual error being \textasciitilde5 years. The histogram
of residual errors roughly follows a normal distribution with mean at
zero so for practical purposes we can assume the errors are normally
distributed.

To further extend our analysis we will incorporate another covariate
from our dataset to our single linear regression model to improve the
performance of the model. We have explored the variables in our dataset
at the start of this section. The variables in general tell us about the
demographics of patients, the types of cancers and treatments, and
various characteristics of the cancers.

We choose the binary categorical variable \texttt{Metastasis} as the
additional covariate for our linear regression model. We fit the model
with \texttt{AgeDiagnosis} and \texttt{Metastasis} as our covriates and
estimate the parameters.

\begin{verbatim}
Residuals:
    Min      1Q  Median      3Q     Max 
-8.8408 -3.1953 -0.4716  2.4082 12.9116 

Coefficients:
              Estimate Std. Error t value Pr(>|t|)    
(Intercept)   13.19498    2.06157   6.400 1.02e-09 ***
AgeDiagnosis   0.89459    0.04535  19.728  < 2e-16 ***
MetastasesYes  1.13560    0.64446   1.762   0.0795 .  
---
Signif. codes:  0 ‘***’ 0.001 ‘**’ 0.01 ‘*’ 0.05 ‘.’ 0.1 ‘ ’ 1

Residual standard error: 4.609 on 207 degrees of freedom
Multiple R-squared:   0.66, Adjusted R-squared:  0.6567 
F-statistic: 200.9 on 2 and 207 DF,  p-value: < 2.2e-16
\end{verbatim}

We observe that, the p-value for the \texttt{AgeDiagnosis} is still
close to zero like our previous regression model. The p-value for our
additional covariate is 0.0795 which is significant at 10\% level. So,
we can reject the null hypothesis at 10\% significance level and say
that the metastases of the cancer has impact on the patient's age at the
death. The estimated residual error is 4.609 which is slightly less than
our previous model(4.632).

The estimated intercept is 13.19 which is slightly less than the the
previous model. The estimated slope is 0.9 for \texttt{AgeDiagnosis} and
1.13 for \texttt{Metastases}. The estimates slope for
\texttt{Metastases} is greater than 1 which indicates that, on average,
patients with metastases have an estimated mean age at death
approximately 1.13 years higher than those without metastases. In
real-world cases, we'd expect patients with metastasized cancer to have
lower age at death. But, patients diagnosed with a cancer could die
before the cancer had metastasized or the cancer is easier to diagnose
once it has metastasized so we get people with higher age at diagnosis
and subsequently higher age at death.

\begin{figure}

{\centering \includegraphics{assignment1_files/figure-latex/reg1-1} 

}

\caption{Fig 4- A regression model for Age at Death conditional on Age at Diagnosis and Metastases.}\label{fig:reg1}
\end{figure}

The linear regression plot shows a positive relationship and a good fit
for our data.The R-squared value is 0.66 which is slightly better than
our previous model(0.655). The histogram of residual errors roughly
follows a normal distribution with mean at zero.

From the plot, we see that there is not much difference between patients
with metastasized cancer and patients with not metastasized cancer. All
the data points are close to each other along the regression line. The
estimated coefficient of the metastasized covariate is also 1.13 (close
to 1), so for patients with same age at diagnosis and different
metastasis cancer state their age at death will only vary by an year.

Overall, we see a slight improvement in our new model which included an
additional covariate alongwith \texttt{AgeDiagnosis}.

We know that, the patients are taking treatment in 12 different
hospitals. To estimate the effect or the non-effect of hospitals on
patient age at death we fit a random intercept model where we consider
\texttt{AgeDiagnosis} to be the fixed effect and cluster the patients by
hospitals by using \texttt{SiteID} as the random effect. We estimate the
parameters of this random intercept model using the \texttt{lmer()}
function in R. To make things easier, we have encoded our
\texttt{SiteID} variable to range from 1 to 12.

\begin{verbatim}
Linear mixed model fit by maximum likelihood  ['lmerMod']
Formula: AgeDeath ~ AgeDiagnosis + (1 | SiteID1)
   Data: CWsample

     AIC      BIC   logLik deviance df.resid 
  1238.6   1251.9   -615.3   1230.6      206 

Scaled residuals: 
     Min       1Q   Median       3Q      Max 
-1.94751 -0.79853 -0.07794  0.55510  2.99798 

Random effects:
 Groups   Name        Variance Std.Dev.
 SiteID1  (Intercept)  2.174   1.475   
 Residual             19.353   4.399   
Number of obs: 210, groups:  SiteID1, 12

Fixed effects:
             Estimate Std. Error t value
(Intercept)  14.06641    2.02125   6.959
AgeDiagnosis  0.89045    0.04355  20.446

Correlation of Fixed Effects:
            (Intr)
AgeDiagnoss -0.964
\end{verbatim}

\textbf{Fixed Effects:} The estimated intercept for the fixed part of
our model i.e.~\texttt{AgeDiagnosis} is 14.07, so it indicates that on
average there is a difference of about 14 years between age at death and
age at diagnosis. The slope or the coefficient of \texttt{AgeDiagnosis}
is 0.89, so on average for each one-unit increase in
\texttt{AgeDiagnosis}, the predicted mean age at death increases by
approximately a year. The t-value is 20.45 which is greater than zero
and hence statistically significant. So, \texttt{AgeDiagnosis} is a
strong predictor of \texttt{AgeDeath}.

\textbf{Random Effects:} The estimated variance between intercepts is
2.17 which isn't very high. It indicates that, the random intercepts are
close to each other and don't vary much. The estimated residual variance
is 19.35 indicating the variability in age at death within the sites.
The higher value of residual variance suggests that there is more
unexplained variability in age at death within the sites, which could be
due to randomness or external factors.

\textbf{Correlation of Fixed Effects:} The estimated correlation value
is -0.964, which indicates a strong negative correlation between the
intercept and \texttt{AgeDiagnosis}. This suggests that as the age at
diagnosis increases, the intercept tends to decrease, and vice versa.

We now calculate the ``Intraclass Correlation Coefficient'' (ICC) for
our model.

\begin{verbatim}
ICC <- between_group_variance / (between_group_variance + residual_variance)
\end{verbatim}

The estimated ICC value is 0.1. It indicates that approximately 10\% of
the total variance in the age at death can be explained by the
differences between the groups or clusters in the data. A value of 0.1
suggests that there is still considerable variability in
\texttt{AgeDeath} not explained by clustering of sites.

\begin{figure}

{\centering \includegraphics{assignment1_files/figure-latex/reg2-1} 

}

\caption{Fig 5- A radom intercept model for Age at Death conditional on Age at Diagnosis and clustered by sites.}\label{fig:reg2}
\end{figure}

We observe that, the intercepts are close to each other as expected from
the estimate of intercept variance. It is difficult to observe any other
significant features from the plot because the random intercepts and the
data points are all close to each other.

The random intercept model makes the assumption that the residual errors
and the random intercepts are normally distributed.It also makes the
assumption of ``Homoscedasticity'' which says that the variability of
the residuals is assumed to be constant across all levels. It also
assumes a linear relationship between the dependent and independent
variables. We plot the QQ Plots of residuals and random intercepts to
test this assumption.

\begin{figure}

{\centering \includegraphics{assignment1_files/figure-latex/re3-1} 

}

\caption{Fig 6- QQ Plots to test the normality assumptions of the random intercept model.}\label{fig:re3}
\end{figure}

We observe that, the qq plot of residuals curves away from the normal
line on the either ends of the graph. The qq plot of random intercepts
are very close to the normal line but don't fall on it. So, the
residuals and random intercepts violate the assumption of normality.
But, for practical purposes we assume the normality of these two as they
closely mimic a normal distribution.

To test for homoscedasticity we plot the residuals against the fitted
values.

\begin{figure}

{\centering \includegraphics{assignment1_files/figure-latex/re4-1} 

}

\caption{Fig 6- Test for Homoscedasticity.}\label{fig:re4}
\end{figure}

We observe that, the spread of the data points is relatively constant
across different levels of the fitted values. So, our model doesn't
violate the homoscedasticity assumption.

One potential limitation of our model is the assumption of independence
of observations within the sites, which may not always be the case in
practice. There could be correlations present among the individuals
within the same site.

We now fit a random intercept model with an additional covariate
\texttt{Metastases} alongwith \texttt{AgeDiagnosis} and estimate its
parameters.

\begin{verbatim}
Linear mixed model fit by maximum likelihood  ['lmerMod']
Formula: AgeDeath ~ AgeDiagnosis + Metastases + (1 | SiteID1)
   Data: CWsample

     AIC      BIC   logLik deviance df.resid 
  1237.3   1254.1   -613.7   1227.3      205 

Scaled residuals: 
     Min       1Q   Median       3Q      Max 
-1.95812 -0.72872 -0.05577  0.57161  2.96306 

Random effects:
 Groups   Name        Variance Std.Dev.
 SiteID1  (Intercept)  2.194   1.481   
 Residual             19.041   4.364   
Number of obs: 210, groups:  SiteID1, 12

Fixed effects:
              Estimate Std. Error t value
(Intercept)    13.7670     2.0132   6.838
AgeDiagnosis    0.8829     0.0434  20.344
MetastasesYes   1.1152     0.6185   1.803

Correlation of Fixed Effects:
            (Intr) AgDgns
AgeDiagnoss -0.948       
MetastassYs -0.084 -0.095
\end{verbatim}

\textbf{Correlation of Fixed Effects:} The estimated correlation value
is -0.95, which indicates a strong negative correlation between the
intercept and \texttt{AgeDiagnosis}. This suggests that as the age at
diagnosis increases, the intercept tends to decrease, and vice versa.
The estimated correlation between the intercept and \texttt{Metastases}
is -0.084 and the correlation between \texttt{AgeDiagnosis} and
\texttt{Metastases} is -0.095. These are small negative values very
close to zero which indicates that there is no significant correlation
present between the intercept and \texttt{Metastases}, and
\texttt{AgeDiagnosis} and \texttt{Metastases}.

\begin{figure}

{\centering \includegraphics{assignment1_files/figure-latex/re575-1} 

}

\caption{Fig 7- Random Intercept Model with additional covariate.}\label{fig:re575}
\end{figure}

We observe that using the estimated values and the plots, this affect of
adding an additional covariate doesn't vary much sites.

We plot the QQ Plots of residuals and random intercepts to test this
assumption.

\begin{figure}

{\centering \includegraphics{assignment1_files/figure-latex/re9-1} 

}

\caption{Fig 8- QQ Plots to test the normality assumptions of the random intercept model.}\label{fig:re9}
\end{figure}

We observe that, the qq plot of residuals curves away from the normal
line on the either ends of the graph. The qq plot of random intercepts
are very close to the normal line but don't fall on it. So, the
residuals and random intercepts violate the assumption of normality.

To test for homoscedasticity we plot the residuals against the fitted
values.

\begin{figure}

{\centering \includegraphics{assignment1_files/figure-latex/re44-1} 

}

\caption{Fig 9- Test for Homoscedasticity.}\label{fig:re44}
\end{figure}

We observe that, the spread of the data points is relatively constant
across different levels of the fitted values. So, our model doesn't
violate the homoscedasticity assumption.

To test whether this random intercept model with addition covariate is
preferable than the single level linear regression model with additional
covariate, we use the Chi-bar test. We obtain the chi-squared value of
7.3842 using \texttt{anova} in R and use this to calculate the p-value
using the following formula in R.

\begin{verbatim}
p =  0.5 * (1 - pchisq(7.3842, df = 1))
\end{verbatim}

The p-value is 0.0033 which is significant at 0.1\% level. We can reject
the null hypotheses at the 0.1\% significance level, in favour of the
hypothesis that the Random Intercept Model is correct or more
preferable.

\hypertarget{conclusion}{%
\subsubsection{Conclusion}\label{conclusion}}

We can conclude the age at diagnosis has a positive impact on the age at
death. The metastatic state of cancer also has a positive effect on the
age at the death. Using these two covariates in the linear regression
model improve the performance of the task. The hospital at which the
patient is taking the treatment doesn't affect the age at death at a
significant level. The cluster by site and using the two covariates in
the random intercept model is the preferable choice for this task. This
analysis didn't produce any real life impact as was expected.

\hypertarget{designing-a-trial}{%
\subsection{Designing a Trial}\label{designing-a-trial}}

\hypertarget{trial-design}{%
\subsubsection{Trial Design}\label{trial-design}}

We propose trial that is a randomized controlled trial (RCT), to
evaluate the effectiveness of interventions. The randomized allocation
of participants to either the intervention or control group helps
minimize selection bias and ensures comparability between groups.

\textbf{Endpoint:} The primary endpoint of the trial is the age at death
of cancer patients. This will be measured as a continuous variable,
aiming to detect whether the intervention increases the expected age at
death by two and a half years or more.

\textbf{Parameters:} We will estimate the parameters required for the
sample size using our dataset. These parameters include the mean age at
death and its standard deviation in the control group, as well as the
expected increase in age at death due to the intervention.

The underlying model for analysis is a linear regression model, where
the age at death is the dependent variable and the intervention group (0
for control, 1 for intervention) is the independent variable. This model
assumes that the relationship between the intervention and age at death
is linear, which might not always be the case. This model will assess
whether there is a statistically significant difference in the age at
death between the intervention and control groups.

\hypertarget{sample-size-requirements}{%
\subsubsection{Sample Size
Requirements}\label{sample-size-requirements}}

To calculate the sample size, we need to specify the type I error rate
(α) and the type II error rate (β). For this analysis, we use the
industry standard values of α = 0.025 (one-sided) and β = 0.2. Using
these standard error rates, we calculate the sample size required to
detect an effect size of two and a half years or more in the age at
death between the intervention and control groups. We do these
calculations in R.

\begin{verbatim}
alpha <- 0.025  # Type I error rate
beta <- 0.2     # Type II error rate
delta <- 2.5   
sigma <- sd(CWsample$AgeDeath)


# Calculate sample size
n <- 2*sigma^2 * (qnorm(1 - alpha) - qnorm(beta))^2 / delta^2
\end{verbatim}

The calculated sample size comes out to be 156 patients per arm.

These sample size calculations assumes that the outcome variable follows
a normal distribution, the outcome variable doesn't have high
variability. As high variable outcome will require a very large sample
size which is not always feasible. So it is assumes that the standard
error of the estimated treatment effect is approximately same as the
standard error of the dependent variable.

\hypertarget{sensitivity-analysis}{%
\subsubsection{Sensitivity Analysis}\label{sensitivity-analysis}}

We now observe how the power of the trial varies according to the
treatment effect and sample size.

\begin{figure}

{\centering \includegraphics{assignment1_files/figure-latex/re24-1} 

}

\caption{Fig 10- Power as a function of Treatment Effect.}\label{fig:re24}
\end{figure}

We observe that we can't attain the required level of power with a very
small size for our treatment effect. We observe that for a large sample
size the power of trial saturates at 1 very quickly which is not
desirable and for a smaller sample size the power doesn't reach 1.

\begin{figure}

{\centering \includegraphics{assignment1_files/figure-latex/re242-1} 

}

\caption{Fig 11- Power as a function of Sample Size.}\label{fig:re242}
\end{figure}

\newpage

\includegraphics[width=10.41667in,height=\textheight]{C:/Athang/MSc DSA/Mixed Models with Medical Applications/aca integrity final f.pdf}

\end{document}
